% ===== wrap-up =====
\begin{frame}{이번 주차 정리}
  \begin{enumerate}
    \item \kb{1.1 RL의 정체성} --- 데이터(에이전트가 만든다) $\cdot$
      신호(정답 $y$가 아니라 보상) $\cdot$ 시간(행동이 미래를 바꾼다)
      $\cdot$ 새 딜레마(탐험 vs 활용). \textbf{리턴}
      $G_t = \sum \gamma^k R_{t+k}$가 최대화 대상이고,
      로봇 시뮬레이션은 연속 상태$\cdot$행동 + 물리 법칙이라
      \textbf{표 기반만으로는 안 됨}.
    \item \kb{1.2 신경망 뼈대} --- 활성 함수 없이는 2층이 선형 하나
      (XOR로 검증), 역전파 = ``오차를 한 층씩 거슬러 올려보낸다''
      (1-2-1 손계산). \textbf{새로움은 학습 알고리즘이 아니라
      손실함수의 정의} --- \texttt{zero\_grad() $\to$ backward()
      $\to$ step()} 세 줄은 지도/강화 \textbf{그대로}.
    \item \kb{1.3 Gymnasium} --- \texttt{reset}/\texttt{step} 계약,
      \texttt{terminated}(진짜 종료) vs \texttt{truncated}(시간
      한도), 시드 두 개를 모두 고정하는 재현성.
      \textbf{무작위 정책 베이스라인 = 평균 24.1스텝} --- 이후
      모든 알고리즘의 ``성공''을 잴 기준.
  \end{enumerate}
  \vskip1em
  \begin{block}{다음 주차 --- 2주차. 멀티암 밴딧}
    상태도, 전이도 없는 \textbf{가장 단순한} RL 문제에서
    ``탐험 vs 활용'' 딜레마를 \textbf{순수한} 형태로 만남.
    $\varepsilon$-greedy와 UCB로 처음 풀어보며 --- CartPole의
    24.1스텝을 넘는 것이 첫 번째 관문.
  \end{block}
\end{frame}
